\section{课堂 Q\&A：偏好、安全、工具与未解的对齐问题}

主 slide 在 30:54 左右结束，但课程一半的信息量来自随后 Q\&A。这里不按问题出现顺序机械抄录，而把讨论重组为五个工程主题：人机分工、uncertainty、preference payload、tool-assisted evaluation，以及 outer/inner alignment 与 interpretability。

\subsection{Machines 做认知劳动，Humans 负责偏好输入}

讲者的长期图景不是让人类亲自阅读所有材料、写所有测试、做所有计算，而是让 AI 承担这些可扩展的 cognitive labor；人类集中处理自己最不可替代的输入---什么结果值得追求、哪些 tradeoff 可接受、何时风险过高。本节先解释最后一张 slide 的责任划分，再用后续 Q\&A 检验这项分工会在哪些地方重新把判断压力推回人类：不确定性、偏好代表性、工具权限和证据信任。

\lecturefigure{slide-17-machines-vs-humans.jpg}{讲座结尾的分工：机器读、算、写代码和核验；人类沟通偏好。}{Stanford Online 官方视频 00:29:45--00:30:54。}

\begin{knowledgebox}{读图：偏好输入仍然需要认知工作}
``Humans: communicate preferences'' 不是只按一个喜欢/不喜欢按钮。人类需要澄清目标、比较冲突价值、决定证据标准、识别不可接受风险，并判断 assistant 是否在操纵评价过程。机器承担更多认知劳动的目的，是把有限人类注意力集中到这些不可简单外包的决策上。
\end{knowledgebox}

\subsection{Uncertainty calibration 与 critique signal-to-noise}

学生先问 hallucination 是否能通过模型不确定性解决。讲者认为教模型在不确定时承认不确定 ``应该是可解的''，但 2023 年仍没有令人满意的状态。简单多次采样也不是完整答案：同一 fine-tuned model 的样本可能高度相关，而独立预训练多个模型成本极高。

\begin{warningbox}{什么是 calibration}
Calibration 要求模型声称 80\% 把握的一组事件大约有 80\% 正确；它不同于语言上加一句 ``我可能错了''。对生成模型，还需定义事件、把长答案拆成可验证 claims，并防止模型用模糊措辞逃避责任。表达不确定性有用，却不能替代事实验证。
\end{warningbox}

Critique 还存在 signal-to-noise 问题。讲者承认很多 critique 只是 nitpicking 或 garbage；在 ``给人类提供检查灵感'' 的高召回模式中，这可能尚可接受，因为人类能快速丢弃无关建议。但如果 assistant 直接做决策，低 precision 会消耗注意力、制造 false alarm，并让真正严重的问题淹没在噪声中。

\subsection{Preference payload：谁的偏好、怎样更新、如何防投毒}

RLHF 经常被画成 ``human feedback''，但 Q\&A 追问了更困难的 payload：谁被视为 human，怎样代表不同群体，偏好冲突如何处理。讲者直言，当时 labeler pool 很大程度上只是公司能够雇到的人，并不多样或具有代表性；有些研究最好由大科技公司之外的机构承担。

\begin{importantbox}{不要直接在任意公开反馈上在线训练}
讲者首先否定了一个危险捷径：把产品界面中的任意用户输入或反馈原样用于持续训练。协调攻击、恶意示范和数据投毒会把反馈渠道变成控制面。安全流程至少需要身份与质量控制、异常检测、分层抽样、审计、保留集和延迟更新，而不是 ``更多反馈自然更民主''。
\end{importantbox}

偏好也会漂移。重新收集 comparisons 并 fine-tune 的计算相对便宜，但模型更新会改变输出风格和行为，依赖旧 prompt 的应用必须重新适配。于是 alignment update 同时是 compatibility event：系统要在更符合新偏好与保持下游稳定之间权衡。

\begin{table}[H]
\centering
\caption{Preference data 的四个边界}
\begin{tabularx}{\textwidth}{L{2.5cm} X X}
\toprule
边界 & 关键问题 & 失败方式 \\
\midrule
Representation & 谁进入 labeler pool？ & 少数群体与专业需求被平均掉 \\
Aggregation & 冲突偏好怎样组合？ & 把规范冲突误写成普通标签噪声 \\
Security & 谁能影响反馈渠道？ & 协调操纵、投毒与奖励劫持 \\
Drift & 偏好与知识怎样随时间更新？ & 行为突变、prompt 兼容性下降 \\
\bottomrule
\end{tabularx}
\end{table}

\subsection{Browsing、API 与 external evidence}

学生讨论了 browsing、calculators、search 与其他 APIs。工具可以让模型引用外部资料、运行计算和检查实时事实，从而把 ``凭参数记忆回答'' 改成 ``给出可复查证据''。但工具同时扩大 action space：模型可能访问不可信来源、泄露数据、调用高权限操作，或者把工具错误包装成确定结论。

\begin{knowledgebox}{工具把 alignment 从文本扩展到系统边界}
只生成 token 时，主要问题是答案质量；能调用工具后，还要定义 credentials、网络域、写权限、预算、日志、回滚和 human approval。Tool result 也不是天然 ground truth：搜索排名、网页内容、API 数据和代码执行环境都可能被操纵。可扩展监督因工具更强，也因工具更危险。
\end{knowledgebox}

在教育场景中，工具还可能造成 overreliance。学生若只接受流畅答案而不练习建模、推导和检验，会把辅助能力误当成自己的理解。好的系统应暴露证据、分解步骤、鼓励用户预测与检查，而不是只输出最终结论。

\subsection{Outer alignment、inner alignment 与 distribution shift}

后段 Q\&A 讨论 inner optimizer。Outer alignment 关注外部训练信号是否真的表达人类目标；inner alignment 关注模型内部学出的优化过程是否在训练分布外继续追随同一目标。讲者的个人判断是，先把 outer signal 做得非常可靠可能尤其重要，因为可信信号还能在新分布上继续训练和校正内部 optimizer。

\[
\underbrace{\text{trusted outer signal}}_{\text{what is rewarded}}
\quad+\quad
\underbrace{\text{coverage under shift}}_{\text{where it remains valid}}
\quad\Longrightarrow\quad
\text{continued correction of learned behavior}.
\]

其中，trusted outer signal 表示人类愿意依赖的评价信号，coverage under shift 表示该信号在新任务与新状态上仍能判断质量，箭头表示一种研究希望而非已证明保证。

\teachervoice{讲者多次使用 ``I think''、``unclear'' 与 ``in practice''。他的论点是：若外部信号足够可信，部分 inner-alignment 风险可以被重新表述为 distribution shift 与持续评价问题；他没有声称 inner alignment 已因此消失。}

\begin{warningbox}{自我 critique 不能凭空创造 ground truth}
让模型审查自己的回答、让两个模型互相辩论或让 reward model 评价 policy，都可能提高覆盖；但如果所有角色共享同一盲点或同一错误目标，循环只会更自洽。必须保留独立证据、对抗样本、可验证任务和人类可理解的 failure signal。
\end{warningbox}

\subsection{Interpretability：有用，但可能 neither sufficient nor necessary}

最后的讨论非常克制。Interpretability 可以帮助发现 deception、追踪模型为何给出某个答案，并作为评价工具箱的一部分；但即使能解释部分内部机制，也不自动告诉我们怎样把这些解释转成可靠训练目标。更尖锐的风险是 selection effect：若只淘汰被当前工具检测出的 misalignment，训练或筛选可能偏向更难检测的失败模式。

\begin{importantbox}{Interpretability 的三层问题}
第一层是 detection：能否看到与失败相关的内部特征；第二层是 attribution：能否说明这些特征如何导致行为；第三层是 control：能否在不破坏其他能力的情况下可靠改变行为。完成第一层不等于完成第三层，因此 interpretability 可能很有用，却既非充分条件，也未必是唯一必要路径。
\end{importantbox}

另一方面，若系统在所有相关分布上的行为都能被可靠评价和约束，内部 ``理由'' 是否必须完全可读仍是开放问题。讲者没有给出确定答案，而是把目标重新落在 ``really, really good evaluation signal''：能否持续判断模型实际做了什么，并据此训练它做我们真正想要的事。

\subsection{本章小结}

Q\&A 把 alignment 从一个 loss 扩展成完整评价系统：需要校准不确定性、控制 critique 噪声、治理 preference payload、防止投毒、管理工具权限、处理 distribution shift，并理解 interpretability 的作用边界。没有任何单项技术能替代 trusted evaluation。

